% MIKU 与 Apollo 路径--速度解耦规划链的融合关系
\begin{tikzpicture}[
    scale=0.88, transform shape,
    input/.style={rectangle, rounded corners=2pt, draw=deepblue,
        fill=inputyellow, minimum width=2.0cm, minimum height=1.05cm,
        align=center, font=\footnotesize},
    miku/.style={rectangle, rounded corners=3pt, draw=deeppink, very thick,
        fill=improvedpink!55, minimum width=4.0cm, minimum height=1.55cm,
        align=center, font=\footnotesize},
    apollo/.style={rectangle, rounded corners=3pt, draw=deepblue,
        fill=lightblue!45, minimum width=2.35cm, minimum height=1.15cm,
        align=center, font=\footnotesize},
    output/.style={rectangle, rounded corners=2pt, draw=deepblue,
        fill=outputpink, minimum width=1.8cm, minimum height=1.05cm,
        align=center, font=\footnotesize},
    mainflow/.style={-Stealth, thick, deepblue},
    mikuflow/.style={-Stealth, very thick, deeppink},
    lab/.style={font=\scriptsize, align=center, fill=white, inner sep=1.5pt},
]

\node[input]  (in)   at (0,0) {参考线、车辆状态\\障碍物预测轨迹};
\node[miku]   (miku) at (3.6,0) {\textbf{MIKU 时变通行带}\\交互感知边界 $+$ 时变投影\\Top-K 空间同伦选择};
\node[apollo] (pqp)  at (7.4,0) {Apollo 路径 QP\\求解 $l(s)$};
\node[apollo] (st)   at (10.3,0) {Apollo ST\\边界映射};
\node[apollo] (sqp)  at (13.2,0) {Apollo 速度 QP\\求解 $s(t)$};
\node[output] (out)  at (16.0,0) {时空轨迹\\$(x,y,v,t)$};

\draw[mainflow] (in) -- (miku);
\draw[mainflow] (miku) -- node[lab, below] {中心路径边界} (pqp);
\draw[mainflow] (pqp) -- node[lab, below] {$l(s)$} (st);
\draw[mainflow] (st) -- node[lab, below] {原 ST 边界} (sqp);
\draw[mainflow] (sqp) -- (out);

\draw[mikuflow] (miku.north) .. controls (5.8,2.05) and (11.0,2.05) .. (sqp.north)
    node[lab, midway, above=1pt, text=deeppink] {\textbf{多冲突点时间同伦 $\mathcal H_t$}\\融合为 ST 位置上下边界};

\node[font=\scriptsize, text=deeppink, anchor=west] at (1.6,-1.35)
    {\rule{0.75cm}{1.1pt}\;MIKU 约束构造与传递};
\node[font=\scriptsize, text=deepblue, anchor=west] at (8.8,-1.35)
    {\rule{0.75cm}{1.1pt}\;Apollo 连续优化主干};

\end{tikzpicture}
